\documentclass[11pt]{article}
\usepackage[margin=1in]{geometry}
\usepackage{amsmath,amssymb,amsthm,enumitem,hyperref}

\newtheorem{theorem}{Theorem}
\theoremstyle{remark}
\newtheorem*{remark}{Remark}

\newcommand{\bigsqcap}{\mathop{\mathchoice{\vcenter{\hbox{\Large$\sqcap$}}}{\sqcap}{\sqcap}{\sqcap}}}

\title{Metamonovalued Reloids Are Monovalued (OPG-57403)}
\author{Dewayne Barnes}
\date{November 26, 2026}

\begin{document}
\maketitle

We give a self-contained, Lean~4--verified characterization settling Conjecture~1023 (\texttt{OPG-57403}~\cite{opg57403}) of Victor Porton~\cite{porton} together with Porton's converse (Theorem~1022).
The core diagonal/antidiagonal argument for $(4) \Rightarrow (1)$ was also given in an October~2026 unrefereed preprint by Ferudun~\cite{ferudun}; this note contributes a proper-meet strengthening $(5) \Rightarrow (1)$ (showing the result does not depend on improper meets $g \sqcap h = \bot$), an explicit $G = \emptyset$ counterexample, and a standalone zero-dependency Lean~4 formalization.

\paragraph{Notation.}
A \emph{reloid} $f : A \to B$ is a (possibly improper) filter on $A \times B$, ordered by reverse inclusion $f \sqsubseteq g \iff g \subseteq f$, with least (improper) element $\bot = \mathcal{P}(A \times B)$, greatest element $\top = \{A \times B\}$, principal reloids $\uparrow R = \{U \subseteq A \times B : R \subseteq U\}$, and meets $\bigsqcap G$ given by the filter generated by $\bigcup_{g \in G} g$ (in particular, $W \in f \sqcap g \iff \exists U \in f,\, V \in g,\, U \cap V \subseteq W$).
The composition $g \circ f : A \to C$ of $f : A \to B$ and $g : B \to C$ is the filter generated by $\{S \circ R : R \in f,\, S \in g\}$, where $S \circ R = \{(x,z) : \exists y \in B,\, (x,y) \in R \wedge (y,z) \in S\}$; the converse is $f^{-1} = \{R^{-1} : R \in f\}$ with $R^{-1} = \{(y,x) : (x,y) \in R\}$.
Write $\Delta_B = \{(b,b) : b \in B\}$ and $\overline{\Delta}_B = (B \times B) \setminus \Delta_B$.

\begin{theorem}
For any reloid $f : A \to B$, the following are equivalent:
\begin{enumerate}[label=(\arabic*)]
  \item $f$ is \emph{monovalued}: there exists a right-unique (partial-functional) relation $R \in f$, i.e., $(x,y_1), (x,y_2) \in R \implies y_1 = y_2$ (equivalently, $f \circ f^{-1} \sqsubseteq \uparrow\Delta_B$);
  \item $f$ is \emph{metamonovalued}: $(\bigsqcap G) \circ f = \bigsqcap_{g \in G} (g \circ f)$ for every set $C$ and every nonempty family $G$ of reloids $B \to C$;
  \item $f$ is \emph{weakly metamonovalued}: $(g \sqcap h) \circ f = (g \circ f) \sqcap (h \circ f)$ for every set $C$ and all reloids $g, h : B \to C$;
  \item $f$ preserves binary meets of principal endoreloids on $B$;
  \item $f$ preserves binary meets $(g \sqcap h) \circ f = (g \circ f) \sqcap (h \circ f)$ for every set $C$ and all reloids $g, h : B \to C$ with $g \sqcap h \neq \bot$.
\end{enumerate}
\end{theorem}

\begin{proof}
First, the parenthetical equivalence in $(1)$ is immediate: if $R \in f$ is right-unique, then $R \circ R^{-1} \subseteq \Delta_B$ gives $\Delta_B \in f \circ f^{-1}$, i.e., $f \circ f^{-1} \sqsubseteq \uparrow\Delta_B$; conversely, $\Delta_B \in f \circ f^{-1}$ yields $R_1, R_2 \in f$ with $R_2 \circ R_1^{-1} \subseteq \Delta_B$, so $R_1 \cap R_2 \in f$ is right-unique.

\medskip\noindent
$(1) \Rightarrow (2)$ (Porton~\cite{porton}, Thm.~1022).
For each $g \in G$, $\bigsqcap G \sqsubseteq g$ gives $(\bigsqcap G) \circ f \sqsubseteq g \circ f$ by monotonicity of composition, so $(\bigsqcap G) \circ f \sqsubseteq \bigsqcap_{g \in G} (g \circ f)$.
Conversely, let $R_0 \in f$ be right-unique and $W \in (\bigsqcap G) \circ f$, so $S \circ R_1 \subseteq W$ for some $R_1 \in f$ and $S \in \bigsqcap G$.
Then $R = R_0 \cap R_1 \in f$ is right-unique and $S \circ R \subseteq W$.
Since $G$ is nonempty (fix $g_0 \in G$), we have $S \supseteq \bigcap_{i=1}^n S_i$ for some $n \ge 1$ and $S_i \in g_i \in G$ (taking $n=1$, $S_1 = B \times C \in g_0$ if needed).
Right-uniqueness of $R$ implies $(\bigcap_{i=1}^n S_i) \circ R = \bigcap_{i=1}^n (S_i \circ R)$: if $(x,z)$ lies in every $S_i \circ R$, the intermediate witnesses $y_i$ with $(x,y_i) \in R$ and $(y_i,z) \in S_i$ must all coincide.
Since each $S_i \circ R \in g_i \circ f \subseteq \bigsqcap_{g \in G} (g \circ f)$, we obtain $W \supseteq \bigcap_{i=1}^n (S_i \circ R) \in \bigsqcap_{g \in G} (g \circ f)$, proving $\bigsqcap_{g \in G} (g \circ f) \sqsubseteq (\bigsqcap G) \circ f$.

\medskip\noindent
$(2) \Rightarrow (3) \Rightarrow (4)$ and $(3) \Rightarrow (5)$ are immediate specializations.

\medskip\noindent
$(4) \Rightarrow (1)$.
Set $g = \uparrow\Delta_B$ and $h = \uparrow\overline{\Delta}_B$ on $B$.
Since $\Delta_B \cap \overline{\Delta}_B = \emptyset$, we have $g \sqcap h = \bot$ and hence $(g \sqcap h) \circ f = \bot \circ f = \bot$ (as $\emptyset \circ (A \times B) = \emptyset$).
By $(4)$, $(g \circ f) \sqcap (h \circ f) = \bot$, so $\emptyset \in (g \circ f) \sqcap (h \circ f)$.
Thus there exist disjoint $U \in g \circ f$ and $V \in h \circ f$, yielding $R_1, R_2 \in f$ with $\Delta_B \circ R_1 \subseteq U$ and $\overline{\Delta}_B \circ R_2 \subseteq V$.
Then $R = R_1 \cap R_2 \in f$ is right-unique: if $(x, y_1), (x, y_2) \in R$ with $y_1 \neq y_2$, then $(x, y_2) \in (\Delta_B \circ R_1) \cap (\overline{\Delta}_B \circ R_2) \subseteq U \cap V = \emptyset$, a contradiction.

\medskip\noindent
$(5) \Rightarrow (1)$.
If $B = \emptyset$, then $A \times B = \emptyset \in f$ is vacuously right-unique.
Otherwise $B \neq \emptyset$; let $C = B \cup \{\ast\}$ with $\ast \notin B$, and define $R_g = (B \times \{\ast\}) \cup \Delta_B$ and $R_h = (B \times \{\ast\}) \cup \overline{\Delta}_B$ in $B \times C$.
Then $R_g \cap R_h = B \times \{\ast\} \neq \emptyset$, so $g = \uparrow R_g$ and $h = \uparrow R_h$ satisfy $g \sqcap h = \uparrow(B \times \{\ast\}) \neq \bot$.
Since $(B \times \{\ast\}) \circ (A \times B) \subseteq A \times \{\ast\}$, we have $A \times \{\ast\} \in (g \sqcap h) \circ f = (g \circ f) \sqcap (h \circ f)$ by $(5)$.
Hence there exist $U \in g \circ f$ and $V \in h \circ f$ with $U \cap V \subseteq A \times \{\ast\}$, giving $R_1, R_2 \in f$ with $R_g \circ R_1 \subseteq U$ and $R_h \circ R_2 \subseteq V$.
Then $R = R_1 \cap R_2 \in f$ is right-unique: if $(x, y_1), (x, y_2) \in R$ with $y_1 \neq y_2$, then $(x, y_2) \in (\Delta_B \circ R_1) \cap (\overline{\Delta}_B \circ R_2) \subseteq (R_g \circ R_1) \cap (R_h \circ R_2) \subseteq U \cap V \subseteq A \times \{\ast\}$, contradicting $y_2 \in B$.
\end{proof}

\begin{remark}
Nonemptiness of $G$ is necessary for $(1) \Rightarrow (2)$: when $G = \emptyset$, $(\bigsqcap \emptyset) \circ f = \top \circ f$ fails to equal $\top$ for any nonempty set $C$ and any non-entire proper monovalued reloid $f$ (i.e., $\exists R \in f$ with $\operatorname{dom}(R) \neq A$)---for example, $f = \uparrow\{(\mathrm{true}, ())\}$ on $\mathrm{Bool} \times \mathrm{Unit}$ with $C = \mathrm{Unit}$ gives $\top \circ f = f \neq \top$.
Conjecture~1023 itself (metamonovalued $\Rightarrow$ monovalued) holds under both the nonempty-$G$ and unrestricted readings.
An accompanying standalone Lean~4 formalization is provided in \texttt{UnsolvedMath/OPG57403.lean}.
\end{remark}

\begin{thebibliography}{3}
\bibitem{porton} V.~Porton, \emph{Algebraic General Topology, Volume~1}, 2013--2019 (Theorem~1022 and Conjecture~1023).
\bibitem{opg57403} UnsolvedMath / Open Problem Garden, \emph{Every metamonovalued reloid is monovalued} (\texttt{OPG-57403}), \url{https://www.unsolvedmath.com/problems/OPG-57403}.
\bibitem{ferudun} A.~Ferudun, \emph{Two Problems of Porton on Reloids: Metamonovalued Reloids Are Monovalued, and $S(S(f)) = S(f)$ Fails}, Zenodo preprint, Oct.~2026, \url{https://zenodo.org/records/23245513}.
\end{thebibliography}

\end{document}
